
\subsection{命题 7.2.35 的证明}
\label{subsec:7-2-5}

将对 $\|(\mathfrak d,\beta)\|$ 归纳，
构造多值截面
$\mathfrak s_{\mathfrak d,k,\beta,\vec P}$。
为确保归纳可行，先证几个引理。

\begin{lemma}{7.2.45}\label{lem:7-2-45}
$\|\beta\|\ge-1$，且
$\|\beta\|=-1$ 当且仅当 $\beta=\beta_0$。
\end{lemma}
\begin{proof}
由定义立即可得。
\end{proof}

\begin{lemma}{7.2.46}\label{lem:7-2-46}
若 $\mathfrak d_\varnothing$ 是唯一的映射
$\varnothing\to\mathbb Z_{\ge0}$，
则 $\|(\mathfrak d_\varnothing,\beta)\|\ge0$。
若 $\mathfrak d\ne\mathfrak d_\varnothing$，
则 $\|(\mathfrak d,\beta)\|>
\max_i\mathfrak d(i)$。
\end{lemma}
\begin{proof}
第二个断言由式~\eqref{eq:7-2-33}
及引理~\ref{lem:7-2-45} 立即得出；
当 $\beta=\beta_0$ 时，
这里要用到稳定性条件给出的 $k\ge2$。
第一个断言由式~\eqref{eq:7-2-34} 得出，
因为此时 $\beta\ne\beta_0$。
\end{proof}

\begin{lemma}{7.2.47}\label{lem:7-2-47}
设
$(\Sigma,w,\mathfrak d)
=(\Sigma_1,w_1,\mathfrak d_1)
\mathbin{\#_\ell}
(\Sigma_2,w_2,\mathfrak d_2^{\max})$
如式~\eqref{eq:7-2-41}，
$\beta_1+\beta_2=\beta$，其中
$(\Sigma_1,w_1,\mathfrak d_1)
\in\mathcal M^{\mathrm{main},\mathfrak d_1}_{k_1+1}(\beta_1)$，
$(\Sigma_2,w_2,\mathfrak d_2^{\max})
\in\mathcal M^{\mathrm{main},\mathfrak d_2^{\max}}_{k_2+1}(\beta_2)$，
而 $(\Sigma,w,\mathfrak d)
\in\mathcal M^{\mathrm{main},\mathfrak d}_{k+1}(\beta)$。
再令 $\mathfrak d_2$ 如式~\eqref{eq:7-2-43}。
则对 $i=1,2$，下列三种情形之一成立：

\noindent\textup{(7.2.48.1)}\quad
$\|(\mathfrak d_i,\beta_i)\|
<\|(\mathfrak d,\beta)\|$；

\noindent\textup{(7.2.48.2)}\quad
$\|(\mathfrak d_i,\beta_i)\|
=\|(\mathfrak d,\beta)\|$，且
$k_i+\|\beta_i\|<k+\|\beta\|$；

\noindent\textup{(7.2.48.3)}\quad
$\|(\mathfrak d_i,\beta_i)\|
=\|(\mathfrak d,\beta)\|$，
$k_i+\|\beta_i\|=k+\|\beta\|$，且
$\|\beta_i\|<\|\beta\|$。
\end{lemma}

\begin{proof}
由于
$\|(\mathfrak d_2^{\max},\beta_2)\|
\ge\|(\mathfrak d_2,\beta_2)\|$，
只须考虑
$\mathfrak d_2=\mathfrak d_2^{\max}$；
以下均作此假设。于是
\begin{equation}
k=k_1+k_2-1,
\tag{7.2.49.1}\label{eq:7-2-49-1}
\end{equation}
\begin{equation}
\|\beta\|
\ge\|\beta_1\|+\|\beta_2\|+1.
\tag{7.2.49.2}\label{eq:7-2-49-2}
\end{equation}
此外有
\begin{equation}
k_i\le1\quad\Longrightarrow\quad\beta_i\ne\beta_0,
\tag{7.2.50.1}\label{eq:7-2-50-1}
\end{equation}
\begin{equation}
k_2\ge1,\qquad k_1\ge0.
\tag{7.2.50.2}\label{eq:7-2-50-2}
\end{equation}

\par\medskip\noindent\begin{minipage}{\textwidth}
\centering
\setlength{\unitlength}{0.8mm}
\begin{picture}(145,64)
\linethickness{0.5pt}
\put(50,31){\circle{42}}
\put(92,31){\circle{42}}
\linethickness{0.25pt}
\qbezier[16](50,57)(24,57)(24,31)
\qbezier[16](24,31)(24,5)(50,5)
\qbezier[10](50,5)(61,5)(68,12)
\qbezier[16](92,57)(118,57)(118,31)
\qbezier[16](118,31)(118,5)(92,5)
\qbezier[10](92,5)(81,5)(74,12)
\linethickness{0.5pt}
\put(71,31){\circle*{1.7}}
\put(50,31){\makebox(0,0){$\beta_1$}}
\put(92,31){\makebox(0,0){$\beta_2$}}
\put(57,49){\makebox(0,0){$\scriptstyle\times$}}
\put(62,47){$\scriptstyle z_\ell$}
\put(80,48){\makebox(0,0){$\scriptstyle\times$}}
\put(77,51){$\scriptstyle z_{\ell-1}$}
\put(104,45){\makebox(0,0){$\scriptstyle\times$}}
\put(107,45){$\scriptstyle z_1$}
\put(113,31){\makebox(0,0){$\scriptstyle\times$}}
\put(117,30){$\scriptstyle z_0$}
\put(108,17){\makebox(0,0){$\scriptstyle\times$}}
\put(111,14){$\scriptstyle z_k$}
\put(76,17){\makebox(0,0){$\scriptstyle\times$}}
\put(77,11){$\scriptstyle z_{\ell+k_1}$}
\put(58,13){\makebox(0,0){$\scriptstyle\times$}}
\put(61,7){$\scriptstyle z_{\ell+k_1-1}$}
\put(33,15){\makebox(0,0){$\scriptstyle\times$}}
\put(29,31){\makebox(0,0){$\scriptstyle\times$}}
\put(50,52){\makebox(0,0){$\scriptstyle\times$}}
\end{picture}
\par{\renewcommand{\thefigure}{7.2.2}\refstepcounter{figure}
\normalfont\bfseries 图~\thefigure\label{fig:7-2-2}}
\end{minipage}\par\medskip

为便于比较，置下列三个最大值：
\[
\begin{aligned}
d_1&=\max\{\mathfrak d_1(i)\mid i=1,\ldots,k_1\},\\
d_2&=\max\{\mathfrak d_2(i)\mid i\ne\ell,
\ i=1,\ldots,k_2\},\\
d&=\max\{\mathfrak d(1),\ldots,\mathfrak d(k)\}.
\end{aligned}
\]
若定义 $d_1$ 的指标集为空（$k_1=0$），
就置 $d_1=0$；
若定义 $d_2$ 的指标集为空（$k_2=1$），
就置 $d_2=0$；
若定义 $d$ 的指标集为空
（$k=0$，等价于 $k_1=0$ 且 $k_2=1$），
就置 $d=0$。
于是 $d=\max\{d_1,d_2\}$。

先证 $i=1$ 的断言。
若 $\beta_2=\beta_0$，则 $k_2\ge2$，从而
\[
\|(\mathfrak d_1,\beta_1)\|
=k_1+\|\beta_1\|+d_1
<k+\|\beta\|+d
=\|(\mathfrak d,\beta)\|.
\]
即 (7.2.48.1) 成立。
若 $\beta_2\ne\beta_0$，则
$\|\beta_1\|<\|\beta\|$ 且 $k_2\ge1$，故
\[
\|(\mathfrak d_1,\beta_1)\|
=k_1+\|\beta_1\|+d_1
<k+\|\beta\|+d
=\|(\mathfrak d,\beta)\|.
\]
此时 (7.2.48.1) 也成立。

再证 $i=2$ 的断言。注意
\[
\mathfrak d_2(\ell)
=\|(\mathfrak d_1,\beta_1)\|
=k_1+\|\beta_1\|+d_1.
\]
分别考虑 $d=d_1$ 和 $d=d_2$。

先设 $d=d_1$。由式~\eqref{eq:7-2-49-1}、
\eqref{eq:7-2-49-2} 得
\[
\begin{aligned}
\|(\mathfrak d,\beta)\|
&=k+\|\beta\|+d\\
&\ge k_1+k_2+\|\beta_1\|+\|\beta_2\|+d_1\\
&\ge k_2+\|\beta_2\|
+\|(\mathfrak d_1,\beta_1)\|.
\end{aligned}
\]
又由引理~\ref{lem:7-2-46}，
\[
\mathfrak d_2(\ell)
=\|(\mathfrak d_1,\beta_1)\|
\ge d_1=d\ge d_2.
\]
因而
$\max_{1\le i\le k_2}\mathfrak d_2(i)
=\|(\mathfrak d_1,\beta_1)\|$，
所以
\begin{equation}
\|(\mathfrak d_2,\beta_2)\|
=k_2+\|\beta_2\|
+\|(\mathfrak d_1,\beta_1)\|
\le\|(\mathfrak d,\beta)\|.
\tag{7.2.51}\label{eq:7-2-51}
\end{equation}
若此处严格小于，则 (7.2.48.1) 成立。
以下设式~\eqref{eq:7-2-51} 取等号。
当 $k_1>1$ 时，
$k=k_1+k_2-1>k_2$，且
$\|\beta\|\ge\|\beta_2\|$，
故 (7.2.48.2) 成立。
当 $k_1=1$ 时，
$\beta_1\ne\beta_0$，从而
$\|\beta\|>\|\beta_2\|$，
又 $k=k_2$，故 (7.2.48.2) 仍成立。
当 $k_1=0$ 时，
$d_1=d=0\le\|(\mathfrak d_1,\beta_1)\|$。
由式~\eqref{eq:7-2-51} 的等号可得
\[
k+\|\beta\|
=\|(\mathfrak d,\beta)\|
=k_2+\|\beta_2\|
+\|(\mathfrak d_1,\beta_1)\|
\ge k_2+\|\beta_2\|.
\]
又 $k_2=k+1>k$，
故 (7.2.48.2) 或 (7.2.48.3) 成立。

再设 $d=d_2>d_1$，分为两种情形：
\noindent\textup{(7.2.52.1)}\quad
$d_2\le\|(\mathfrak d_1,\beta_1)\|$；

\noindent\textup{(7.2.52.2)}\quad
$d_2>\|(\mathfrak d_1,\beta_1)\|$。

在 (7.2.52.1) 中，
记 $A=\|(\mathfrak d_1,\beta_1)\|$。
由 $d_2\le A=\mathfrak d_2(\ell)$，
有
$\|(\mathfrak d_2,\beta_2)\|
=k_2+\|\beta_2\|+A$。
另一方面，由式~\eqref{eq:7-2-49-1}、
\eqref{eq:7-2-49-2}，
\[
\|(\mathfrak d,\beta)\|
\ge k_1+k_2+\|\beta_1\|+\|\beta_2\|+d.
\]
若 $k_1>0$，则
$A=k_1+\|\beta_1\|+d_1$，
两式之差至少为 $d-d_1>0$；
若 $k_1=0$，则 $d_1=0$、
$A=\|\beta_1\|$，
两式之差至少为 $d>0$。
故 (7.2.48.1) 成立。
\translatornote{原书在情形 (7.2.52.1)
把 $\|(\mathfrak d_1,\beta_1)\|$
先等同于 $d_2$，又在后续等式链中
把它等同于 $d_1$；这些等号均不由该情形推出。
此处直接依 (7.2.33)、(7.2.49) 与
$d_2\le\|(\mathfrak d_1,\beta_1)\|$
比较，所得严格不等式仍是原欲证结论。
参见疑误 A-0436。}

在 (7.2.52.2) 中，
\[
\begin{aligned}
\|(\mathfrak d,\beta)\|
&=k+\|\beta\|+d\\
&\ge k_1+k_2+\|\beta_1\|+\|\beta_2\|+d_2\\
&\ge k_2+\|\beta_2\|+d_2
=\|(\mathfrak d_2,\beta_2)\|.
\end{aligned}
\]
这里用到了 $k_1+\|\beta_1\|\ge0$，
它由稳定性条件给出。
若取等号，则 $k_1+\|\beta_1\|=0$，
所以 $\beta_1\ne\beta_0$ 且 $k_1=0$。
于是 $k=k_2-1$、
$\|\beta\|>\|\beta_2\|$，
故 (7.2.48.2) 或 (7.2.48.3) 成立。
引理证毕。
\end{proof}

现在开始证明命题~\ref{prop:7-2-35}。
固定两个正数 $\delta$ 与 $K$。
对每个 $\|(\mathfrak d,\beta)\|\le K$，
构造多值截面系
$\mathfrak s_{\mathfrak d,k,\beta,\vec P}$。
构造按归纳进行。为组织归纳步骤，
在四元组 $(\mathfrak d,k,\beta,\vec P)$
的集合上引入次序 $<$：
当且仅当下列条件之一成立时，称
$(\mathfrak d',k',\beta',\vec P')
<(\mathfrak d,k,\beta,\vec P)$：

\noindent\textup{(7.2.53.1)}\quad
$\|(\mathfrak d',\beta')\|
<\|(\mathfrak d,\beta)\|$；

\noindent\textup{(7.2.53.2)}\quad
$\|(\mathfrak d',\beta')\|
=\|(\mathfrak d,\beta)\|$，且
$k'+\|\beta'\|<k+\|\beta\|$；

\noindent\textup{(7.2.53.3)}\quad
$\|(\mathfrak d',\beta')\|
=\|(\mathfrak d,\beta)\|$、
$k'+\|\beta'\|=k+\|\beta\|$，且
$\|\beta'\|<\|\beta\|$；

\noindent\textup{(7.2.53.4)}\quad
$\|(\mathfrak d',\beta')\|
=\|(\mathfrak d,\beta)\|$、
$k'+\|\beta'\|=k+\|\beta\|$、
$\|\beta'\|=\|\beta\|$，且
$\sum_i\deg P_i'>\sum_i\deg P_i$。

将归纳地证明如下两步：

\noindent\textup{(7.2.54.1)}\quad
假设所有
$(\mathfrak d',k',\beta',\vec P')
<(\mathfrak d,k,\beta,\vec P)$
对应的多值截面
$\mathfrak s_{\mathfrak d',k',\beta',\vec P'}$
都已定义，
且 $g<\|(\mathfrak d,\beta)\|$
时的 $\mathcal X_g(L)$ 已定义，
则定义满足所需性质的
$\mathfrak s_{\mathfrak d,k,\beta,\vec P}$。

\noindent\textup{(7.2.54.2)}\quad
假设 $g=\|(\mathfrak d,\beta)\|$
时的全部
$\mathfrak s_{\mathfrak d,k,\beta,\vec P}$
均已定义，则定义 $\mathcal X_g(L)$。

先处理 $\|(\mathfrak d,\beta)\|=0$：
同时定义 $\mathcal X_0(L)$
及相应的多值截面。
引理~\ref{lem:7-2-46} 表明，
此时 $\mathfrak d=\mathfrak d_\varnothing$、
$\|\beta\|=0$；特别地 $\beta\ne\beta_0$。
考虑模空间
$\mathcal M^{\mathrm{main},\mathfrak d}_{0+1}(\beta)$。
由于 $\|\beta\|=0$，
它不能冒泡，故无边界；
其相容条件因而为空。
于是可自由选取其上的多值截面
$\mathfrak s_{\mathfrak d_\varnothing,0,\beta,\vec 1}$，
并使它与零截面横截。
这里 $\vec 1$ 对应唯一的映射
$\varnothing\to S(L)$。
取所有形如
\[
\operatorname{ev}_{0*}
\bigl(\mathfrak s^{-1}_{\mathfrak d_\varnothing,
0,\beta,\vec 1}(0)\bigr)
\qquad(\|\beta\|=0)
\]
的链，选择其光滑三角剖分，
并考虑剖分中出现的全部奇异单形。
按注~\ref{rem:7-2-15}(3) 的方法，
把每个单形识别为奇异单形；
以后相似情形均使用此法，不再重复说明。
若其中某单形已出现过，便将其舍去。
所得集合即 $\mathcal X_0(L)$，
这完成归纳的第一步。

接着说明 (7.2.54.1) 的归纳步骤。
引理~\ref{lem:7-2-46} 给出
$\mathfrak d(i)<\|(\mathfrak d,\beta)\|$，
故要求 $P_i\in\mathcal X_{\mathfrak d(i)}(L)$
在此时已有意义。
考察相容条件~\ref{cond:7-2-38}
与~\ref{cond:7-2-44}。
设
$(\mathbf p_m,(p_{1,m},\ldots,p_{k,m}))$
是
$\mathcal M^{\mathrm{main},\mathfrak d}_{k+1}
(L;\beta;\vec P)$
的仓西邻域内收敛于边界点的序列。
如前所述，此极限落在
相容条件~\ref{cond:7-2-38}
或~\ref{cond:7-2-44} 的情形中。
引理~\ref{lem:7-2-47} 因而表明，
极限处
$\mathfrak s_{\mathfrak d,k,\beta,\vec P}$
的值已经由归纳假设确定。
这里还必须用到下述引理。

\begin{lemma}{7.2.55}\label{lem:7-2-55}
相容条件~\ref{cond:7-2-38}
与~\ref{cond:7-2-44} 确定的
$\mathfrak s_{\mathfrak d,k,\beta,\vec P}$
边界值，在两种情形重叠的点彼此一致。
\end{lemma}

暂且假设引理~\ref{lem:7-2-55}，
先完成命题~\ref{prop:7-2-35} 的证明。
应用 \textbf{[FuOn99II]} 的定理 3.11
（即附录 A1 的定理 A1.23），
可以延拓边界值，得到与零截面横截的
$\mathfrak s_{\mathfrak d,k,\beta,\vec P}$。
其零集由引理 A1.26 可作三角剖分；
可以使此剖分与先前归纳步骤
给出的边界剖分相合。
令 $g=\|(\mathfrak d,\beta)\|$。
取零集剖分中的全部单形，
视为光滑奇异单形；
舍去已属于 $\mathcal X_{(g-1)}(L)$
或先前步骤已选出的单形。
其余单形定为第 $g$ 代，
收入 $\mathcal X_g(L)$。
如此可选得该子集，
使 $\mathcal X_{(g)}(L)$ 对此零集
具有所需的生成性质。
由此完成 (7.2.54.1)。

当 $\dim P_1=\dim P_2=0$、
$\beta=\beta_0$ 时，
$\mathcal M_{2+1}(\beta_0;\vec P)$
没有边界，相容条件为空；
因此也可自由选取其上的
$\mathfrak s_{\mathfrak d,2,\beta_0,\vec P}$。

步骤 (7.2.54.2) 较为容易。
把 (7.2.54.1) 中所得的
所有第 $g$ 代的链取并；
除 (7.2.28.2) 外，它已具有全部所需性质。
再适当地加入可数多个奇异单形，
使 (7.2.28.2) 成立，
同时保持 (7.2.28.3)；
参见引理~\ref{lem:7-2-12} 证明的最后一步。

最后检查上述构造能否在给定仓西邻域内完成，
即是否满足 (7.2.36.5)。
这正是第~7.2.3 小节提到的问题。
对所有
$\|(\mathfrak d,\beta)\|\le K$
的多值截面，归纳步数有上界 $C(K)$；
注意满足该界的
$\beta\in\mathfrak G(L)$ 只有有限个。
起初选取一个与原截面
（即仓西结构定义中的仓西映射）
相距至多 $\epsilon^{C(K)}\delta$
的多值截面；
$\epsilon$ 是稍后选取的充分小正数。
把归纳步骤安排为还满足：

\noindent\textup{(7.2.56.1)}\quad
假设第 $I$ 次归纳步骤所得的多值截面系
与原仓西映射相距至多
$\epsilon^{C(K)-I}\delta$，
其零集与原仓西映射的零集
也相距至多
$\epsilon^{C(K)-I}\delta$。

\noindent\textup{(7.2.56.2)}\quad
在第 $I+1$ 步构造
$\mathfrak s_{\mathfrak d,k,\beta,\vec P}$ 时，
取第 $I$ 步及以前的多值截面的纤维积，
从而得到边界
$\partial\mathcal M^{\mathrm{main},\mathfrak d}_{k+1}
(\beta;\vec P)$ 上的多值截面
$\mathfrak s^{\partial}_{\mathfrak d,k,\beta,\vec P}$。

\noindent\textup{(7.2.56.3)}\quad
由构造可知，
$\mathfrak s^{\partial}_{\mathfrak d,k,\beta,\vec P}$
与原仓西映射相距至多
$C\epsilon^{C(K)-I}\delta$，
其零点集与原仓西映射的零集
也相距至多
$C\epsilon^{C(K)-I}\delta$。
这里 $C$ 可由所用模空间及 $K$ 控制。

\noindent\textup{(7.2.56.4)}\quad
故可取延拓的多值截面
$\mathfrak s_{\mathfrak d,k,\beta,\vec P,\epsilon}$，
使它与仓西映射相距至多
$C^2\epsilon^{C(K)-I}\delta$，
其零集与原仓西映射的零集
也相距至多
$C^2\epsilon^{C(K)-I}\delta$。
这同样由构造得出。

\noindent\textup{(7.2.56.5)}\quad
选取 $C^2\epsilon<1$，
则 (7.2.56.4) 保证归纳可继续进行。
再把 $\delta$ 取得足够小，
便可使扰动后仓西映射的零集
远离仓西邻域的端部；
因此第~7.2.3 小节的问题不会出现。

一般还须处理无穷多个光滑奇异单形 $P_i$。
但障碍丛是从伪全纯圆盘模空间上的障碍丛拉回的，
仓西邻域则由这些单形与模空间仓西邻域
的横截纤维积给出。
因此，在 $\|(\mathfrak d,\beta)\|\le K$ 的约束下，
只须研究有限多个模空间，
即可使扰动模空间远离仓西邻域边界。
由此可把常数 $C$ 等的选取
做得与奇异单形 $P_i$ 无关。
式组 (7.2.56) 也给出 (7.2.36.5)，
命题~\ref{prop:7-2-35} 证毕。

\begin{remark}{7.2.57}\label{rem:7-2-57}
只要构造扰动多值截面仅需有限多个步骤，
式组 (7.2.56) 的论证便可用于更一般的情形：
令扰动多值截面充分接近原仓西映射，
使其零集不触及仓西邻域的端部。
因此后面的小节不再重复此论证。
\end{remark}

\begin{proof}[引理~\ref{lem:7-2-55} 的证明]
两条相容条件在重叠处有三种组合：
第一种须两次使用相容条件~\ref{cond:7-2-38}；
第二种各使用一次相容条件~\ref{cond:7-2-38}
与~\ref{cond:7-2-44}；
第三种须两次使用相容条件~\ref{cond:7-2-44}。
这里只讨论第三种；
其余两种较易，略去。

为简化记号，只考虑如下情形，
一般情形相似。
$\Sigma$ 的边界上有两个奇点
$p_0,p'_0$，
图~\ref{fig:7-2-3} 显示两种分解。

\par\medskip\noindent\begin{minipage}{\textwidth}
\centering
\setlength{\unitlength}{0.75mm}
\begin{picture}(145,90)
\linethickness{0.5pt}
\put(31,67){\circle{24}}
\put(55,67){\circle{24}}
\put(79,67){\circle{24}}
\put(43,67){\circle*{1.5}}
\put(67,67){\circle*{1.5}}
\put(31,67){\makebox(0,0){$\Sigma_4$}}
\put(55,67){\makebox(0,0){$\Sigma_3$}}
\put(79,67){\makebox(0,0){$\Sigma_2$}}
\put(40,82){$\scriptstyle p'_0$}
\put(66,82){$\scriptstyle p_0$}
\put(21.9,74.9){\line(1,1){1.2}}\put(21.9,76.1){\line(1,-1){1.2}}
\put(38.9,74.9){\line(1,1){1.2}}\put(38.9,76.1){\line(1,-1){1.2}}
\put(21.9,57.9){\line(1,1){1.2}}\put(21.9,59.1){\line(1,-1){1.2}}
\put(36.4,56){\line(1,1){1.2}}\put(36.4,57.2){\line(1,-1){1.2}}
\put(45.9,74.9){\line(1,1){1.2}}\put(45.9,76.1){\line(1,-1){1.2}}
\put(54.4,78.4){\line(1,1){1.2}}\put(54.4,79.6){\line(1,-1){1.2}}
\put(60.9,76.4){\line(1,1){1.2}}\put(60.9,77.6){\line(1,-1){1.2}}
\put(62.9,57.9){\line(1,1){1.2}}\put(62.9,59.1){\line(1,-1){1.2}}
\put(69.9,74.9){\line(1,1){1.2}}\put(69.9,76.1){\line(1,-1){1.2}}
\put(78.4,78.4){\line(1,1){1.2}}\put(78.4,79.6){\line(1,-1){1.2}}
\put(86.9,74.9){\line(1,1){1.2}}\put(86.9,76.1){\line(1,-1){1.2}}
\put(91,67){\circle*{1.5}}
\put(86.9,57.9){\line(1,1){1.2}}\put(86.9,59.1){\line(1,-1){1.2}}
\put(69.9,57.9){\line(1,1){1.2}}\put(69.9,59.1){\line(1,-1){1.2}}
\put(39,77.5){$\scriptstyle z_{\ell+\ell'-1}$}
\put(45,80.5){$\scriptstyle z_{\ell+\ell'-2}$}
\put(63,78){$\scriptstyle z_\ell$}
\put(30,52.5){$\scriptstyle z_{\ell+\ell'+k_4-2}$}
\put(55,52.5){$\scriptstyle z_{\ell+k_4+k_3-2}$}
\put(72,71){$\scriptstyle z_{\ell-1}$}
\put(83,71){$\scriptstyle z_1$}
\put(94,66.5){$\scriptstyle z_0$}
\put(83.5,60.5){$\scriptstyle z_k$}
\put(78,52.5){$\scriptstyle z_{\ell+k_4+k_3-1}$}
\put(31,23){\circle{24}}
\put(55,23){\circle{24}}
\put(79,23){\circle{24}}
\put(43,23){\circle*{1.5}}
\put(67,23){\circle*{1.5}}
\put(31,23){\makebox(0,0){$\Sigma_4$}}
\put(55,23){\makebox(0,0){$\Sigma_3$}}
\put(79,23){\makebox(0,0){$\Sigma_2$}}
\put(40,38){$\scriptstyle p'_0$}
\put(66,38){$\scriptstyle p_0$}
\put(21.9,30.9){\line(1,1){1.2}}\put(21.9,32.1){\line(1,-1){1.2}}
\put(38.9,30.9){\line(1,1){1.2}}\put(38.9,32.1){\line(1,-1){1.2}}
\put(21.9,13.9){\line(1,1){1.2}}\put(21.9,15.1){\line(1,-1){1.2}}
\put(36.4,12){\line(1,1){1.2}}\put(36.4,13.2){\line(1,-1){1.2}}
\put(45.9,30.9){\line(1,1){1.2}}\put(45.9,32.1){\line(1,-1){1.2}}
\put(54.4,34.4){\line(1,1){1.2}}\put(54.4,35.6){\line(1,-1){1.2}}
\put(60.9,32.4){\line(1,1){1.2}}\put(60.9,33.6){\line(1,-1){1.2}}
\put(62.9,13.9){\line(1,1){1.2}}\put(62.9,15.1){\line(1,-1){1.2}}
\put(46.9,13.9){\line(1,1){1.2}}\put(46.9,15.1){\line(1,-1){1.2}}
\put(69.9,30.9){\line(1,1){1.2}}\put(69.9,32.1){\line(1,-1){1.2}}
\put(78.4,34.4){\line(1,1){1.2}}\put(78.4,35.6){\line(1,-1){1.2}}
\put(86.9,30.9){\line(1,1){1.2}}\put(86.9,32.1){\line(1,-1){1.2}}
\put(91,23){\circle*{1.5}}
\put(86.9,13.9){\line(1,1){1.2}}\put(86.9,15.1){\line(1,-1){1.2}}
\put(69.9,13.9){\line(1,1){1.2}}\put(69.9,15.1){\line(1,-1){1.2}}
\put(39,33.5){$\scriptstyle z_{\ell+\ell'-1}$}
\put(45,36.5){$\scriptstyle z_{\ell+\ell'-2}$}
\put(63,34){$\scriptstyle z_\ell$}
\put(49,8.5){$\scriptstyle z_{\ell+\ell'+k_4+k_3-2}$}
\put(72,27){$\scriptstyle z_{\ell-1}$}
\put(83,27){$\scriptstyle z_1$}
\put(94,22.5){$\scriptstyle z_0$}
\put(83.5,16.5){$\scriptstyle z_k$}
\put(76.5,8.5){$\scriptstyle z_{\ell+\ell'+k_4+k_3-1}$}
\end{picture}
\par{\renewcommand{\thefigure}{7.2.3}\refstepcounter{figure}
\normalfont\bfseries 图~\thefigure\label{fig:7-2-3}}
\end{minipage}\par\medskip

假设在 $p_0$ 处分开 $\Sigma$ 得
$\Sigma_1\cup\Sigma_2$，
其中 $\Sigma_2$ 包含第 $0$ 个（末个）标记点，
而 $p'_0$ 位于 $\Sigma_1$。
另一方面，在 $p'_0$ 处分开 $\Sigma_1$
得 $\Sigma_3\cup\Sigma_4$，
其中 $p_0$ 位于 $\Sigma_3$。
设 $p'_0$ 是 $\Sigma_3$ 的第 $\ell'$ 个标记点，
$p_0$ 是 $\Sigma_2$ 的第 $\ell$ 个标记点。
于是记
\[
\Sigma=(\Sigma_4\mathbin{\#_{\ell'}}\Sigma_3)
\mathbin{\#_\ell}\Sigma_2.
\]

比较两种在边界上获得多值截面的分解：
一为
\[
\Sigma=(\Sigma_4\mathbin{\#_{\ell'}}\Sigma_3)
\mathbin{\#_\ell}\Sigma_2,
\]
另一为
\[
\Sigma=\Sigma_4
\mathbin{\#_{\ell+\ell'-1}}
(\Sigma_3\mathbin{\#_\ell}\Sigma_2).
\]
设 $\Sigma_i$ 有 $k_i+1$ 个边界标记点。
则 $\Sigma$ 有 $k+1$ 个，其中
$k=k_2+k_3+k_4-2$，
而 $\Sigma_1=\Sigma_4\mathbin{\#_{\ell'}}\Sigma_3$
有 $k_1+1$ 个，其中 $k_1=k_3+k_4-1$。
给 $\Sigma_i$ 配上
$\beta_i\in\pi_2(M,L)$，于是
$\beta_1=\beta_3+\beta_4$、
$\beta=\beta_2+\beta_3+\beta_4$。
再给 $\Sigma$ 配上装饰 $\mathfrak d$；
各 $\Sigma_i$ 的装饰按下列方式确定：
\[
\mathfrak d_1(i)=\mathfrak d(i+\ell-1),
\qquad
\mathfrak d_4(i)=\mathfrak d(i+\ell+\ell'-2),
\]
\[
\mathfrak d_2(i)=
\begin{cases}
\mathfrak d(i),&i=1,\ldots,\ell-1,\\
g_1,&i=\ell,\\
\mathfrak d(i+k_1-1),&i=\ell+1,\ldots,k_2,
\end{cases}
\]
\[
\mathfrak d_3(i)=
\begin{cases}
\mathfrak d_1(i),&i=1,\ldots,\ell'-1,\\
g_4,&i=\ell',\\
\mathfrak d_1(i+k_4-1),&i=\ell'+1,\ldots,k_3.
\end{cases}
\]
其中 $g_1\le\|(\mathfrak d_1,\beta_1)\|$、
$g_4\le\|(\mathfrak d_4,\beta_4)\|$
稍后确定。
设 $w\colon\Sigma\to M$ 的类为 $\beta$，
其在 $\Sigma_i$ 上的限制
$w_i\colon\Sigma_i\to M$ 的类为 $\beta_i$。
写作
\begin{equation}
(\Sigma,w,\mathfrak d)
=\bigl(
(\Sigma_4,w_4,\mathfrak d_4)
\mathbin{\#_{\ell'}}
(\Sigma_3,w_3,\mathfrak d_3)
\bigr)
\mathbin{\#_\ell}
(\Sigma_2,w_2,\mathfrak d_2).
\tag{7.2.58}\label{eq:7-2-58}
\end{equation}

对 $P_i\in\mathcal X_{\mathfrak d(i)}(L)$，
先置
\[
\vec P_{(4)}
=(P_{\ell+\ell'-1},\ldots,
P_{\ell+\ell'+k_4-2}),
\]
\[
Q_4=\operatorname{ev}_{0*}\left(
\mathcal M^{\mathrm{main},\mathfrak d_4}_{k_4+1}
(\beta_4;\vec P_{(4)})
^{\mathfrak s_{\mathfrak d_4,k_4,\beta_4,\vec P_{(4)}}}
\right).
\]
由归纳假设，$Q_4$ 可分解为
有限多个代不超过
$\|(\mathfrak d_4,\beta_4)\|$
的奇异单形。
为使记号简洁，
用 $Q_4$ 代替其中一个单形，
记其代为 $g_4$。
接着置
\[
\begin{aligned}
\vec P_{(3)}
={}&(P_\ell,\ldots,P_{\ell+\ell'-2},Q_4,
P_{\ell+\ell'+k_4-1},\ldots,
P_{\ell+k_3+k_4-2}),\\
Q_3={}&\operatorname{ev}_{0*}\left(
\mathcal M^{\mathrm{main},\mathfrak d_3}_{k_3+1}
(\beta_3;\vec P_{(3)})
^{\mathfrak s_{\mathfrak d_3,k_3,\beta_3,\vec P_{(3)}}}
\right).
\end{aligned}
\]
再置
\[
\begin{aligned}
\vec P_{(1)}
&=(P_\ell,\ldots,P_{\ell+k_3+k_4-2}),\\
Q_1&=\operatorname{ev}_{0*}\left(
\mathcal M^{\mathrm{main},\mathfrak d_1}_{k_1+1}
(\beta_1;\vec P_{(1)})
^{\mathfrak s_{\mathfrak d_1,k_1,\beta_1,\vec P_{(1)}}}
\right).
\end{aligned}
\]
注意
$\ell+k_3+k_4-2=\ell+k_1-1$。
图~\ref{fig:7-2-4}、\ref{fig:7-2-5}
分别表示逐个分支与先合并分支的链输出。

\par\medskip\noindent\begin{minipage}{\textwidth}
\centering
\setlength{\unitlength}{0.75mm}
\begin{picture}(145,65)
\linethickness{0.5pt}
\put(28,38){\circle{30}}
\put(58,38){\circle{30}}
\put(88,38){\circle{30}}
\put(28,38){\makebox(0,0){$\Sigma_4$}}
\put(58,38){\makebox(0,0){$\Sigma_3$}}
\put(88,38){\makebox(0,0){$\Sigma_2$}}
\put(43,38){\circle*{1.4}}
\put(73,38){\circle*{1.4}}
\put(18,60){$\scriptstyle\vec P_{(4)}$}
\put(52,60){$\scriptstyle\vec P_{(3)}$}
\put(14.4,44.9){\line(1,1){1.2}}\put(14.4,46.1){\line(1,-1){1.2}}
\put(27.4,52.4){\line(1,1){1.2}}\put(27.4,53.6){\line(1,-1){1.2}}
\put(34.9,50.4){\line(1,1){1.2}}\put(34.9,51.6){\line(1,-1){1.2}}
\put(12.4,37.4){\line(1,1){1.2}}\put(12.4,38.6){\line(1,-1){1.2}}
\put(18.3,25.5){\line(1,1){1.2}}\put(18.3,26.7){\line(1,-1){1.2}}
\put(27.4,22.4){\line(1,1){1.2}}\put(27.4,23.6){\line(1,-1){1.2}}
\put(33.1,23.5){\line(1,1){1.2}}\put(33.1,24.7){\line(1,-1){1.2}}
\put(46.9,44.9){\line(1,1){1.2}}\put(46.9,46.1){\line(1,-1){1.2}}
\put(57.4,52.4){\line(1,1){1.2}}\put(57.4,53.6){\line(1,-1){1.2}}
\put(64.9,50.4){\line(1,1){1.2}}\put(64.9,51.6){\line(1,-1){1.2}}
\put(68,26.8){\line(1,1){1.2}}\put(68,28){\line(1,-1){1.2}}
\put(57.4,22.4){\line(1,1){1.2}}\put(57.4,23.6){\line(1,-1){1.2}}
\put(46.9,26.8){\line(1,1){1.2}}\put(46.9,28){\line(1,-1){1.2}}
\put(74.4,44.9){\line(1,1){1.2}}\put(74.4,46.1){\line(1,-1){1.2}}
\put(87.4,52.4){\line(1,1){1.2}}\put(87.4,53.6){\line(1,-1){1.2}}
\put(94.9,50.4){\line(1,1){1.2}}\put(94.9,51.6){\line(1,-1){1.2}}
\put(98,26.8){\line(1,1){1.2}}\put(98,28){\line(1,-1){1.2}}
\put(76.8,26.8){\line(1,1){1.2}}\put(76.8,28){\line(1,-1){1.2}}
\put(103,38){\circle*{1.4}}
\put(26.5,55){$\scriptstyle P$}
\put(9,49){$\scriptstyle P$}
\put(5,38.5){$\scriptstyle P$}
\put(10,25.5){$\scriptstyle P$}
\put(21,18.5){$\scriptstyle P$}
\put(38,52){$\scriptstyle P_{\ell+\ell'-1}$}
\put(28,19.5){$\scriptstyle P_{\ell+\ell'+k_4-2}$}
\put(43.5,42){$\scriptstyle P_{\ell+\ell'-2}$}
\put(60,55){$\scriptstyle P$}
\put(67,52.5){$\scriptstyle P_\ell$}
\put(48,19.5){$\scriptstyle P_{\ell+\ell'+k_4-1}$}
\put(66,20){$\scriptstyle P$}
\put(68.6,27.4){\line(3,-4){10}}
\put(75,16.5){$\scriptstyle P_{\ell+k_4+k_3-2}$}
\put(72.5,54.5){$\scriptstyle P_{\ell-1}$}
\put(101,49.5){$\scriptstyle P_1$}
\put(101,25.5){$\scriptstyle P_{k-1}$}
\put(70.5,21.5){$\scriptstyle P$}
\put(103,38){\vector(1,0){9}}
\put(114,37){$\scriptstyle z_0$}
\put(43,18){\line(0,-1){5}}
\put(43,13){\vector(-1,0){15}}
\put(33,15){$\scriptstyle Q_4$}
\put(73,18){\line(0,-1){9}}
\put(73,9){\vector(-1,0){25}}
\put(57,5.5){$\scriptstyle Q_3$}
\put(55,15){$\scriptstyle\vec P_{(3)}$}
\end{picture}
\par{\renewcommand{\thefigure}{7.2.4}\refstepcounter{figure}
\normalfont\bfseries 图~\thefigure\label{fig:7-2-4}}
\end{minipage}\par\medskip

\par\medskip\noindent\begin{minipage}{\textwidth}
\centering
\setlength{\unitlength}{0.75mm}
\begin{picture}(145,63)
\linethickness{0.5pt}
\put(28,38){\circle{30}}
\put(58,38){\circle{30}}
\put(88,38){\circle{30}}
\put(28,38){\makebox(0,0){$\Sigma_4$}}
\put(58,38){\makebox(0,0){$\Sigma_3$}}
\put(88,38){\makebox(0,0){$\Sigma_2$}}
\put(43,38){\circle*{1.4}}
\put(73,38){\circle*{1.4}}
\put(38,59.5){$\scriptstyle\vec P_{(1)}$}
\put(4,36){\rotatebox{90}{$\scriptstyle\vec P_{(1)}$}}
\put(14.4,44.9){\line(1,1){1.2}}\put(14.4,46.1){\line(1,-1){1.2}}
\put(34.9,50.4){\line(1,1){1.2}}\put(34.9,51.6){\line(1,-1){1.2}}
\put(12.4,37.4){\line(1,1){1.2}}\put(12.4,38.6){\line(1,-1){1.2}}
\put(16.8,26.8){\line(1,1){1.2}}\put(16.8,28){\line(1,-1){1.2}}
\put(31.3,22.9){\line(1,1){1.2}}\put(31.3,24.1){\line(1,-1){1.2}}
\put(46.9,44.9){\line(1,1){1.2}}\put(46.9,46.1){\line(1,-1){1.2}}
\put(64.9,50.4){\line(1,1){1.2}}\put(64.9,51.6){\line(1,-1){1.2}}
\put(68,26.8){\line(1,1){1.2}}\put(68,28){\line(1,-1){1.2}}
\put(59.4,22.5){\line(1,1){1.2}}\put(59.4,23.7){\line(1,-1){1.2}}
\put(46.8,26.8){\line(1,1){1.2}}\put(46.8,28){\line(1,-1){1.2}}
\put(74.4,44.9){\line(1,1){1.2}}\put(74.4,46.1){\line(1,-1){1.2}}
\put(87.4,52.4){\line(1,1){1.2}}\put(87.4,53.6){\line(1,-1){1.2}}
\put(94.9,50.4){\line(1,1){1.2}}\put(94.9,51.6){\line(1,-1){1.2}}
\put(98,26.8){\line(1,1){1.2}}\put(98,28){\line(1,-1){1.2}}
\put(76.8,26.8){\line(1,1){1.2}}\put(76.8,28){\line(1,-1){1.2}}
\put(103,38){\circle*{1.4}}
\put(9,49){$\scriptstyle P$}
\put(31,54.5){$\scriptstyle P_{\ell+\ell'-1}$}
\put(10,25.5){$\scriptstyle P$}
\put(28,19.5){$\scriptstyle P_{\ell+\ell'+k_4-2}$}
\put(47,54.5){$\scriptstyle P_{\ell+\ell'-2}$}
\put(67,52.5){$\scriptstyle P_\ell$}
\put(48,19.5){$\scriptstyle P_{\ell+\ell'+k_4-1}$}
\put(66,20){$\scriptstyle P$}
\put(68.6,27.4){\line(3,-4){10}}
\put(75,16.5){$\scriptstyle P_{\ell+k_4+k_3-2}$}
\put(72.5,54.5){$\scriptstyle P_{\ell-1}$}
\put(101,49.5){$\scriptstyle P_1$}
\put(101,25.5){$\scriptstyle P_{k-1}$}
\put(70.5,21.5){$\scriptstyle P$}
\put(103,38){\vector(1,0){9}}
\put(114,37){$\scriptstyle z_0$}
\put(73,18){\line(0,-1){8}}
\put(73,10){\vector(-1,0){25}}
\put(57,6.5){$\scriptstyle Q_1$}
\put(38,12){$\scriptstyle\vec P_{(1)}$}
\end{picture}
\par{\renewcommand{\thefigure}{7.2.5}\refstepcounter{figure}
\normalfont\bfseries 图~\thefigure\label{fig:7-2-5}}
\end{minipage}\par\medskip

$Q_1$ 可分解为
$\mathcal X_{(\|(\mathfrak d_1,\beta_1)\|)}(L)$
中的单形；
为简洁起见，令 $Q_1$ 表示其中一个，
其代记为 $g_1$。
于是置
\[
\begin{aligned}
\vec P_{(2)}
&=(P_1,\ldots,P_{\ell-1},Q_1,
P_{\ell+k_3+k_4-1},\ldots,P_k),\\
\vec P'_{(2)}
&=(P_1,\ldots,P_{\ell-1},Q_3,
P_{\ell+k_3+k_4-1},\ldots,P_k).
\end{aligned}
\]

如第~7.2.4 小节，考虑
$\mathcal M^{\mathrm{main},\mathfrak d}_{k+1}
(\beta;P_1,\ldots,P_k)$
的仓西邻域中的点列
$(\mathbf p_m,(p_{1,m},\ldots,p_{k,m}))$。
设 $\mathbf p_\infty=(\Sigma,w,\mathfrak d)$，
并取 $p_{i,\infty}\in P_i$，
使
\[
(\mathbf p_\infty,(p_{1,\infty},\ldots,p_{k,\infty}))
=(\mathbf p_\infty,\vec p_\infty)
\]
仍在该仓西邻域中。
依式~\eqref{eq:7-2-58} 的分解，
两次应用相容条件~\ref{cond:7-2-44}，
得到
\begin{equation}
\begin{aligned}
\mathfrak s_{\mathfrak d,k,\beta,\vec P}
(\mathbf p_\infty,\vec p_\infty)
={}&\mathfrak s_{\mathfrak d_4,k_4,\beta_4,\vec P_{(4)}}
(\mathbf p_{(4),\infty},\vec p_{(4),\infty})\\
&\oplus
\mathfrak s_{\mathfrak d_3,k_3,\beta_3,\vec P_{(3)}}
(\mathbf p_{(3),\infty},\vec p_{(3),\infty})\\
&\oplus
\mathfrak s_{\mathfrak d_2,k_2,\beta_2,\vec P_{(2)}}
(\mathbf p_{(2),\infty},\vec p_{(2),\infty}).
\end{aligned}
\tag{7.2.59}\label{eq:7-2-59}
\end{equation}
这里
$\mathbf p_{(i),\infty}=(\Sigma_i,w_i,\mathfrak d_i)$，
$\vec p_{(i),\infty}
=(w_i(z_1),\ldots,w_i(z_{k_i}))$，
其中 $z_j$ 为边界标记点。

现在用另一方式分解 $(\Sigma,w,\mathfrak d)$：
\[
(\Sigma,w,\mathfrak d)
=(\Sigma_4,w_4,\mathfrak d_4)
\mathbin{\#_{\ell+\ell'-1}}
\bigl((\Sigma_3,w_3,\mathfrak d_3)
\mathbin{\#_\ell}
(\Sigma_2,w_2,\mathfrak d'_2)\bigr),
\]
其中
\[
\mathfrak d'_2(i)=
\begin{cases}
\mathfrak d(i),&i=1,\ldots,\ell-1,\\
g_3,&i=\ell,\\
\mathfrak d(i+k_1-1),&i=\ell+1,\ldots,k_2,
\end{cases}
\]
$g_3\le\|(\mathfrak d_3,\beta_3)\|$
是 $Q_3$ 的代。于是得到
\begin{equation}
\begin{aligned}
\mathfrak s_{\mathfrak d,k,\beta,\vec P}
(\mathbf p_\infty,\vec p_\infty)
={}&\mathfrak s_{\mathfrak d_4,k_4,\beta_4,\vec P_{(4)}}
(\mathbf p_{(4),\infty},\vec p_{(4),\infty})\\
&\oplus
\mathfrak s_{\mathfrak d_3,k_3,\beta_3,\vec P_{(3)}}
(\mathbf p_{(3),\infty},\vec p_{(3),\infty})\\
&\oplus
\mathfrak s_{\mathfrak d'_2,k_2,\beta_2,\vec P'_{(2)}}
(\mathbf p_{(2),\infty},\vec p'_{(2),\infty}).
\end{aligned}
\tag{7.2.60}\label{eq:7-2-60}
\end{equation}
式~\eqref{eq:7-2-60} 与
\eqref{eq:7-2-59} 的差别在第三个直和项：
前者使用带撇的
$\mathfrak d'_2$、$\vec P'_{(2)}$，
后者使用不带撇的
$\mathfrak d_2$、$\vec P_{(2)}$。
式~\eqref{eq:7-2-59} 之所以如此，
是因为在那里把
$(\Sigma_4,w_4,\mathfrak d_4)
\mathbin{\#_{\ell'}}
(\Sigma_3,w_3,\mathfrak d_3)$
视为
$\mathcal M^{\mathrm{main},\mathfrak d_1}_{k_1+1}
(L;\beta_1;\vec P_{(1)})$
的一个仓西邻域中的元素。

此处引理~\ref{lem:7-2-55} 所要证明的，
就是式~\eqref{eq:7-2-59} 与
\eqref{eq:7-2-60} 的右端相等。
回忆 $Q_3$ 位于 $Q_1$ 的边界，
因此是 $\partial Q_1$ 的分解中的一项。
由作为归纳假设的
相容条件~\ref{cond:7-2-38}，
两个多值截面在此重合，
从而得到所需性质。
引理证毕。
\end{proof}
